\documentclass[../../main.tex]{subfiles}% 注意这里的文件路径不能用 ./main.tex ,否则用latexmk编译子文件会报错
\graphicspath{{\subfix{./image/}}{\subfix{../../image/}}} % 指定图片引用路径,后续可以直接使用图片文件名
% 如果图片存放在上级目录,则使用 ../../image/ 作为备用路径

% 例如:
% \begin{figure}[H]
% \centering
% \includegraphics[width=0.3\textwidth]{图.png}
% \caption{}
% \label{figure:图}
% \end{figure}
% 注意:上述\label{}一定要放在\caption{}之后,否则引用图片序号会只会显示??.

\begin{document}

\section{含参量积分}

\subsection{含参量正常积分}

\begin{definition}[含参量积分]
设 \( f(x,y) \) 是定义在矩形区域 \( R = [a,b] \times [c,d] \) 上的二元函数。当 \( x \) 取 \( [a,b] \) 上某定值时，函数 \( f(x,y) \) 则是定义在 \( [c,d] \) 上以 \( y \) 为自变量的一元函数。倘若这时 \( f(x,y) \) 在 \( [c,d] \) 上可积，则其积分值是 \( x \) 在 \( [a,b] \) 上取值的函数，记它为 \( \varphi(x) \)，就有
\begin{align}
\varphi(x) = \int_{c}^{d} f(x,y) \, \mathrm{d}y, \, x \in [a,b]. \label{eq::::::--23u80u11}
\end{align}
一般地，设 \( f(x,y) \) 为定义在区域 \( G = \{ (x,y) \mid c(x) \leqslant y \leqslant d(x), a \leqslant x \leqslant b \} \) 上的二元函数，其中 \( c(x), d(x) \) 为定义在 \( [a,b] \) 上的连续函数，若对于 \( [a,b] \) 上每一固定的 \( x \) 值，\( f(x,y) \) 作为 \( y \) 的函数在闭区间 \( [c(x),d(x)] \) 上可积，则其积分值是 \( x \) 在 \( [a,b] \) 上取值的函数，记作 \( F(x) \) 时，就有
\begin{align}
F(x) = \int_{c(x)}^{d(x)} f(x,y) \, \mathrm{d}y, \, x \in [a,b]. \label{eq::::::--23u80u12}
\end{align}
用积分形式所定义的这两个函数 \(\eqref{eq::::::--23u80u11}\) 与 \(\eqref{eq::::::--23u80u12}\)，通称为定义在 \( [a,b] \) 上\textbf{含参量 \( x \) 的（正常）积分}，或简称\textbf{含参量积分}。
\end{definition}

\begin{theorem}[连续性]\label{theorem:含参量积分的连续性}
若二元函数 \( f(x,y) \) 在矩形区域 \( R = [a,b] \times [c,d] \) 上连续，则函数
\[
\varphi(x) = \int_{c}^{d} f(x,y) \, \mathrm{d}y,\quad \psi(y) = \int_{a}^{b} f(x,y) \, \mathrm{d}x .
\]
都在 \( [a,b] \) 上连续。
\end{theorem}
\begin{remark}
对于这个定理的结论也可以写成如下的形式：若 \( f(x,y) \) 在矩形区域 \( R \) 上连续，则对任何 \( x_0 \in [a,b] \)，都有
\[
\lim_{x \to x_0} \int_{c}^{d} f(x,y) \, \mathrm{d}y = \int_{c}^{d} \lim_{x \to x_0} f(x,y) \, \mathrm{d}y.
\]
这个结论表明，定义在矩形区域上的连续函数，其极限运算与积分运算的顺序是可以交换的。
\end{remark}
\begin{proof}
设 \( x \in [a,b] \)，对充分小的 \( \Delta x \)，有 \( x + \Delta x \in [a,b] \)（若 \( x \) 为区间的端点，则仅考虑 \( \Delta x > 0 \) 或 \( \Delta x < 0 \)），于是
\begin{align}
\varphi(x + \Delta x) - \varphi(x) = \int_{c}^{d} [f(x + \Delta x,y) - f(x,y)] \, \mathrm{d}y. \label{eq:::--------------jiopewjg3}
\end{align}
由于 \( f(x,y) \) 在有界闭域 \( R \) 上连续，从而一致连续，即对任给的正数 \( \varepsilon \)，总存在某个正数 \( \delta \)，对 \( R \) 内任意两点 \( (x_1,y_1) \) 与 \( (x_2,y_2) \)，只要
\[
|x_1 - x_2| < \delta, \, |y_1 - y_2| < \delta,
\]
就有
\begin{align}
|f(x_1,y_1) - f(x_2,y_2)| < \varepsilon.\label{eq:::--------------jiopewjg4}
\end{align}
所以由 \(\eqref{eq:::--------------jiopewjg3}\)，\(\eqref{eq:::--------------jiopewjg4}\) 可推得：当 \( |\Delta x| < \delta \) 时，
\[
|\varphi(x + \Delta x) - \varphi(x)| \leqslant \int_{c}^{d} |f(x + \Delta x,y) - f(x,y)| \, \mathrm{d}y
< \int_{c}^{d} \varepsilon \, \mathrm{d}x = \varepsilon(d - c).
\]
这就证明了 \( \varphi(x) \) 在 \( [a,b] \) 上连续。

同理可证：若 \( f(x,y) \) 在矩形区域 \( R \) 上连续，则含参量 \( y \) 的积分
\[
\psi(y) = \int_{a}^{b} f(x,y) \, \mathrm{d}x .
\]
在 \( [c,d] \) 上连续。
\end{proof}

\begin{theorem}[连续性]\label{theorem:含参量积分的连续性1}
设二元函数 \( f(x,y) \) 在区域
\[
G = \{ (x,y) \mid c(x) \leqslant y \leqslant d(x), a \leqslant x \leqslant b \}
\]
上连续，其中 \( c(x), d(x) \) 为 \( [a,b] \) 上的连续函数，则函数
\begin{align}
F(x) = \int_{c(x)}^{d(x)} f(x,y) \, \mathrm{d}y. \label{eq:::--------------jiopewjg6}
\end{align}
在 \( [a,b] \) 上连续。
\end{theorem}
\begin{proof}
对积分 \(\eqref{eq:::--------------jiopewjg6}\) 用换元积分法，令
\[
y = c(x) + t(d(x) - c(x)).
\]
当 \( y \) 在 \( c(x) \) 与 \( d(x) \) 之间取值时，\( t \) 在 \( [0,1] \) 上取值，且
\[
\mathrm{d}y = (d(x) - c(x)) \, \mathrm{d}t.
\]
所以从 \(\eqref{eq:::--------------jiopewjg6}\) 式可得
\begin{align*}
F(x)  \int_{c(x)}^{d(x)} f(x,y) \, \mathrm{d}y
= \int_{0}^{1} f(x,c(x) + t(d(x) - c(x))) (d(x) - c(x)) \, \mathrm{d}t.
\end{align*}
由于被积函数
\[
f(x,c(x) + t(d(x) - c(x))) (d(x) - c(x))
\]
在矩形区域 \( [a,b] \times [0,1] \) 上连续，由\refthe{theorem:含参量积分的连续性}得积分 \(\eqref{eq:::--------------jiopewjg6}\) 所确定的函数 \( F(x) \) 在 \( [a,b] \) 上连续。
\end{proof}

\begin{theorem}[可微性]\label{theorem:含参量积分的可微性1}
若函数 \( f(x,y) \) 与其偏导数 \( \frac{\partial}{\partial x}f(x,y) \) 都在矩形区域 \( R = [a,b] \times [c,d] \) 上连续，则
\[
\varphi(x) = \int_{c}^{d} f(x,y) \, \mathrm{d}y
\]
在 \( [a,b] \) 上可微，且
\[
\frac{\mathrm{d}}{\mathrm{d}x} \int_{c}^{d} f(x,y) \, \mathrm{d}y = \int_{c}^{d} \frac{\partial}{\partial x} f(x,y) \, \mathrm{d}y.
\]
\end{theorem}
\begin{proof}
对于 \( [a,b] \) 内任一点 \( x \)，设 \( x + \Delta x \in [a,b] \)（若 \( x \) 为区间端点，则讨论单侧导数），则
\[
\frac{\varphi(x + \Delta x) - \varphi(x)}{\Delta x} = \int_{c}^{d} \frac{f(x + \Delta x,y) - f(x,y)}{\Delta x} \, \mathrm{d}y.
\]
由微分学的拉格朗日中值定理及 \( f_x(x,y) \) 在有界闭域 \( R \) 上连续（从而一致连续），对任给正数 \( \varepsilon \)，存在正数 \( \delta \)，只要当 \( |\Delta x| < \delta \) 时，就有
\[
\left| \frac{f(x + \Delta x,y) - f(x,y)}{\Delta x} - f_x(x,y) \right|
= |f_x(x + \theta \Delta x,y) - f_x(x,y)| < \varepsilon,
\]
其中 \( \theta \in (0,1) \)。因此
\begin{align*}
\left| \frac{\Delta \varphi}{\Delta x} - \int_{c}^{d} f_x(x,y) \, \mathrm{d}y \right|
\leqslant \int_{c}^{d} \left| \frac{f(x + \Delta x,y) - f(x,y)}{\Delta x} - f_x(x,y) \right| \, \mathrm{d}y
< \varepsilon(d - c).
\end{align*}
这就证得对一切 \( x \in [a,b] \)，有
\[
\frac{\mathrm{d}}{\mathrm{d}x} \varphi(x) = \int_{c}^{d} \frac{\partial}{\partial x} f(x,y) \, \mathrm{d}y.
\]
\end{proof}

\begin{theorem}[可微性]\label{theorem:含参量积分的可微性2}
设 \( f(x,y), f_x(x,y) \) 在 \( R = [a,b] \times [p,q] \) 上连续，\( c(x), d(x) \) 为定义在 \( [a,b] \) 上其值含于 \( [p,q] \) 内的可微函数，则函数
\[
F(x) = \int_{c(x)}^{d(x)} f(x,y) \, \mathrm{d}y
\]
在 \( [a,b] \) 上可微，且
\begin{align}
F'(x) = \int_{c(x)}^{d(x)} f_x(x,y) \, \mathrm{d}y + f(x,d(x)) d'(x) - f(x,c(x)) c'(x). \label{eq::-------6486-7}
\end{align}
\end{theorem} 
\begin{proof}
把 \( F(x) \) 看作复合函数
\[
F(x) = H(x,c,d) = \int_{c}^{d} f(x,y) \, \mathrm{d}y,
c = c(x), \, d = d(x).
\]
由复合函数求导法则及变限积分的求导法则，有
\[
\frac{\mathrm{d}}{\mathrm{d}x} F(x) = \frac{\partial H}{\partial x} + \frac{\partial H}{\partial c} \frac{\mathrm{d}c}{\mathrm{d}x} + \frac{\partial H}{\partial d} \frac{\mathrm{d}d}{\mathrm{d}x}
= \int_{c(x)}^{d(x)} f_x(x,y) \, \mathrm{d}y + f(x,d(x)) d'(x) - f(x,c(x)) c'(x).
\]
\end{proof}

\begin{theorem}[可积性]\label{theorem:含参量正常积分的可积性定理}
若 \( f(x,y) \) 在矩形区域 \( R = [a,b] \times [c,d] \) 上连续，则 \( \varphi(x) \) 和 \( \psi(y) \) 分别在 \( [a,b] \) 和 \( [c,d] \) 上可积。
这就是说：在 \( f(x,y) \) 连续性假设下，同时存在两个求积顺序不同的积分：
\[
\int_{a}^{b} \left[ \int_{c}^{d} f(x,y) \, \mathrm{d}y \right] \, \mathrm{d}x \quad \text{与} \quad \int_{c}^{d} \left[ \int_{a}^{b} f(x,y) \, \mathrm{d}x \right] \, \mathrm{d}y.
\]
为书写简便起见，今后将上述两个积分写作
\[
\int_{a}^{b} \mathrm{d}x \int_{c}^{d} f(x,y) \, \mathrm{d}y \quad \text{与} \quad \int_{c}^{d} \mathrm{d}y \int_{a}^{b} f(x,y) \, \mathrm{d}x,
\]
前者表示 \( f(x,y) \) 先对 \( y \) 求积然后对 \( x \) 求积，后者则求积顺序相反。它们统称为\textbf{累次积分}，或更确切地称为\textbf{二次积分}。
\end{theorem}
\begin{proof}

\end{proof}

\begin{theorem}
若 \( f(x,y) \) 在矩形区域 \( R = [a,b] \times [c,d] \) 上连续，则
\begin{align}
\int_{a}^{b} \mathrm{d}x \int_{c}^{d} f(x,y) \, \mathrm{d}y = \int_{c}^{d} \mathrm{d}y \int_{a}^{b} f(x,y) \, \mathrm{d}x. \label{eq:::::::----4618-8}
\end{align}
\end{theorem}
\begin{note}
这个定理指出，在 \( f(x,y) \) 连续性假设下，累次积分与求积顺序无关。
\end{note}
\begin{proof}
记
\[
\varphi_1(u) = \int_{a}^{u} \mathrm{d}x \int_{c}^{d} f(x,y) \, \mathrm{d}y,
\varphi_2(u) = \int_{c}^{d} \mathrm{d}y \int_{a}^{u} f(x,y) \, \mathrm{d}x,
\]
其中 \( u \in [a,b] \)，现在分别求 \( \varphi_1(u) \) 与 \( \varphi_2(u) \) 的导数。
\[
\varphi_1'(u) = \frac{\mathrm{d}}{\mathrm{d}u} \int_{a}^{u} \varphi(x) \, \mathrm{d}x = \varphi(u).
\]
对于 \( \varphi_2(u) \)，令 \( H(u,y) = \int_{a}^{u} f(x,y) \, \mathrm{d}x \)，则有
\[
\varphi_2(u) = \int_{c}^{d} H(u,y) \, \mathrm{d}y.
\]
因为 \( H(u,y) \) 与 \( H_u(u,y) = f(u,y) \) 都在 \( R \) 上连续，由\refthe{theorem:含参量积分的可微性1},
\[
\varphi_2'(u) = \frac{\mathrm{d}}{\mathrm{d}u} \int_{c}^{d} H(u,y) \, \mathrm{d}y = \int_{c}^{d} H_u(u,y) \, \mathrm{d}y
= \int_{c}^{d} f(u,y) \, \mathrm{d}y = \varphi(u).
\]
故得 \( \varphi_1'(u) = \varphi_2'(u) \)，因此对一切 \( u \in [a,b] \)，有
\[
\varphi_1(u) = \varphi_2(u) + k \quad (k \text{ 为常数}).
\]
当 \( u = a \) 时，\( \varphi_1(a) = \varphi_2(a) = 0 \)，于是 \( k = 0 \)，即得
\[
\varphi_1(u) = \varphi_2(u), \, u \in [a,b].
\]
取 \( u = b \)，就得到所要证明的 \(\eqref{eq:::::::----4618-8}\) 式。
\end{proof}




\subsection{含参量反常积分}

\subsubsection{含参量反常积分的一致收敛性及其判别法}

\begin{definition}[含参量反常积分]
设函数$f(x,y)$定义在无界区域$R=\{(x,y)|x\in I,c\leqslant y<+\infty\}$上，其中$I$为一区间，若对每一个固定的$x\in I$，反常积分
\begin{align}
\int_{c}^{+\infty}f(x,y)\mathrm{d}y \label{eq::::--wefw234fsd3480t34t3532jrwdr24t}
\end{align}
都收敛，则它的值是$x$在$I$上取值的函数，当记这个函数为$\varPhi(x)$时，则有
\begin{align}
\varPhi(x)=\int_{c}^{+\infty}f(x,y)\mathrm{d}y,\ x\in I, \label{eq::::--wefw234fsd3480t34t32jrwdr24t}
\end{align}
称\eqref{eq::::--wefw234fsd3480t34t3532jrwdr24t}式为定义在$I$上的含参量$x$的\textbf{无穷限反常积分}，或简称\textbf{含参量反常积分}.
\end{definition}

\begin{definition}
若含参量反常积分\eqref{eq::::--wefw234fsd3480t34t3532jrwdr24t}与函数$\varPhi(x)$对任给的正数$\varepsilon$,总存在某一实数$N>c$,使得当$M>N$时,对一切$x\in I$,都有
\begin{align*}
\left| \int_{c}^{M}f(x,y)\mathrm{d}y - \varPhi(x) \right| < \varepsilon,
\end{align*}
即
\begin{align*}
\left| \int_{M}^{+\infty}f(x,y)\mathrm{d}y \right| < \varepsilon,
\end{align*}
则称含参量反常积分\eqref{eq::::--wefw234fsd3480t34t3532jrwdr24t}在$I$上一致收敛于$\varPhi(x)$,或简单地说含参量反常积分\eqref{eq::::--wefw234fsd3480t34t3532jrwdr24t}在$I$\textbf{上一致收敛}.

若对任意$[a,b]\subset I$,含参量反常积分\eqref{eq::::--wefw234fsd3480t34t3532jrwdr24t}在$[a,b]$上一致收敛,则称\eqref{eq::::--wefw234fsd3480t34t3532jrwdr24t}在$I$上\textbf{内闭一致收敛}(若$I=[a,b]$,则内闭一致收敛即一致收敛).
\end{definition}

\begin{theorem}[一致收敛的Cauchy准则]
含参量反常积分\eqref{eq::::--wefw234fsd3480t34t3532jrwdr24t}在$I$上一致收敛的充要条件是:对任给正数$\varepsilon$,总存在某一实数$M>c$,使得当$A_1,A_2>M$时,对一切$x\in I$,都有
\begin{align*}
\left| \int_{A_1}^{A_2}f(x,y)\mathrm{d}y \right| < \varepsilon. 
\end{align*}

\end{theorem}
\begin{proof}

\end{proof}

\begin{theorem}
含参量反常积分$\int_{c}^{+\infty}f(x,y)\mathrm{d}y$在$I$上一致收敛的充分必要条件是
\begin{align*}
\lim_{A \to +\infty}F(A)=0,
\end{align*}
其中$F(A)=\sup_{x\in I}\left| \int_{A}^{+\infty}f(x,y)\mathrm{d}y \right|$.
\end{theorem}
\begin{proof}

\end{proof}

\begin{theorem}\label{theorem:数学分析--定理19.9}
含参量反常积分\eqref{eq::::--wefw234fsd3480t34t3532jrwdr24t}在$I$上一致收敛的充要条件是:对任一趋于$+\infty$的递增数列$\{A_n\}$(其中$A_1 = c$),函数项级数
\begin{align}
\sum_{n=1}^{\infty} \int_{A_n}^{A_{n+1}} f(x,y)\mathrm{d}y = \sum_{n=1}^{\infty} u_n(x) \label{eq::::--wefw234fsd34802jrw4t}
\end{align}
在$I$上一致收敛.
\end{theorem}
\begin{proof}

{\heiti 必要性} 由\eqref{eq::::--wefw234fsd3480t34t3532jrwdr24t}在$I$上一致收敛,故对任给$\varepsilon > 0$,必存在$M > c$,使当$A'' > A' > M$时,对一切$x \in I$,总有
\begin{align}
\left| \int_{A'}^{A''} f(x,y)\mathrm{d}y \right| < \varepsilon. \label{eq::::--wefw234fsd34802jrw4f34t}
\end{align}
又由$A_n \to +\infty\ (n \to \infty)$,所以对正数$M$,存在正整数$N$,只要当$m > n > N$时,就有$A_m \geqslant A_n > M$.由\eqref{eq::::--wefw234fsd34802jrw4f34t}对一切$x \in I$,就有
\begin{align*}
|u_n(x) + \cdots + u_m(x)| &= \left| \int_{A_m}^{A_{m+1}} f(x,y)\mathrm{d}y + \cdots + \int_{A_n}^{A_{n+1}} f(x,y)\mathrm{d}y \right| \\
&= \left| \int_{A_n}^{A_{m+1}} f(x,y)\mathrm{d}y \right| < \varepsilon.
\end{align*}
这就证明了级数\eqref{eq::::--wefw234fsd34802jrw4t}在$I$上一致收敛.

{\heiti 充分性:}用反证法.假若\eqref{eq::::--wefw234fsd3480t34t3532jrwdr24t}在$I$上不一致收敛,则存在某个正数$\varepsilon_0$,使得对于任何实数$M > c$,存在相应的$A'' > A' > M$和$x' \in I$,使得
\begin{align*}
\left| \int_{A'}^{A''} f(x',y)\mathrm{d}y \right| \geqslant \varepsilon_0.
\end{align*}
现取$M_1 = \max\{1,c\}$,则存在$A_2 > A_1 > M_1$及$x_1 \in I$,使得
\begin{align*}
\left| \int_{A_1}^{A_2} f(x_1,y)\mathrm{d}y \right| \geqslant \varepsilon_0.
\end{align*}
一般地,取$M_n = \max\{n,A_{2(n-1)}\}\ (n \geqslant 2)$,则有$A_{2n} > A_{2n-1} > M_n$及$x_n \in I$,使得
\begin{align}
\left| \int_{A_{2n-1}}^{A_{2n}} f(x_n,y)\mathrm{d}y \right| \geqslant \varepsilon_0. \label{eq::::--wefw234fsd34802jrwdr24t}
\end{align}
由上述所得到的数列$\{A_n\}$是递增数列,且$\lim\limits_{n \to \infty} A_n = +\infty$.现在考察级数
\begin{align*}
\sum_{n=1}^{\infty} u_n(x) = \sum_{n=1}^{\infty} \int_{A_n}^{A_{n+1}} f(x,y)\mathrm{d}y.
\end{align*}
由\eqref{eq::::--wefw234fsd34802jrwdr24t}式知存在正数$\varepsilon_0$,对任何正整数$N$,只要$n > N$,就有某个$x_n \in I$,使得
\begin{align*}
|u_{2n-1}(x_n)| = \left| \int_{A_{2n-1}}^{A_{2n}} f(x_n,y)\mathrm{d}y \right| \geqslant \varepsilon_0.
\end{align*}
这与级数\eqref{eq::::--wefw234fsd34802jrw4t}在$I$上一致收敛的假设矛盾.故含参量反常积分\eqref{eq::::--wefw234fsd3480t34t3532jrwdr24t}在$I$上一致收敛.
\end{proof}

\begin{theorem}[含参量反常积分一致收敛判别法]\label{theorem:含参量反常积分一致收敛判别法}
\begin{enumerate}
\item\label{theorem:含参积分Weierstrass判别法(M判别法)} \textbf{Weierstrass判别法(M判别法)} 

设有函数$g(y)$,使得
\[
|f(x,y)| \leqslant g(y),\ (x,y) \in I \times [c, +\infty).
\]
若$\int_{c}^{+\infty} g(y)\mathrm{d}y$收敛,则$\int_{c}^{+\infty} f(x,y)\mathrm{d}y$在$I$上一致收敛.

\item \textbf{Dirichlet判别法} 

设

(i) 对一切实数$N > c$,含参量正常积分
\[
\int_{c}^{N} f(x,y)\mathrm{d}y
\]
对参量$x$在$I$上一致有界,即存在正数$M$,对一切$N > c$及一切$x \in I$,都有
\[
\left| \int_{c}^{N} f(x,y)\mathrm{d}y \right| \leqslant M.
\]

(ii) 对每一个$x \in I$,函数$g(x,y)$为$y$的单调函数,且当$y \to +\infty$时,对参量$x$,$g(x,y)$一致地收敛于0.

则含参量反常积分
\[
\int_{c}^{+\infty} f(x,y)g(x,y)\mathrm{d}y
\]
在$I$上一致收敛.

\item \textbf{Abel判别法} 

设

(i) $\int_{c}^{+\infty} f(x,y)\mathrm{d}y$在$I$上一致收敛.

(ii) 对每一个$x \in I$,函数$g(x,y)$为$y$的单调函数,且对参量$x$,$g(x,y)$在$I$上一致有界.

则含参量反常积分
\[
\int_{c}^{+\infty} f(x,y)g(x,y)\mathrm{d}y
\]
在$I$上一致收敛.
\end{enumerate}
\end{theorem}

\begin{example}
证明:$\int_0^{+\infty}{\frac{\alpha t}{1+\alpha ^2+t^2}e^{-\alpha ^2t^2}\cos \left( \alpha ^2t^2 \right) \,\mathrm{d}t}$在$\alpha \in \left( 0,+\infty \right)$上一致收敛.
\end{example}
\begin{proof}

由于 \(1+\alpha^2+t^2 \geqslant 1+t^2\) 且 \(|\cos(\alpha^2 t^2)| \leqslant 1\)，又求导易知$\underset{x\in \left[ 0,+\infty \right)}{\max}\left\{ xe^{-x^2} \right\} =\frac{1}{\sqrt{2e}}.$故
\begin{align*}
\left| \int_0^{+\infty}{\frac{\alpha t}{1+\alpha ^2+t^2}e^{-\alpha ^2t^2}\cos\mathrm{(}\alpha ^2t^2)\mathrm{d}t} \right|\leqslant \int_0^{+\infty}{\frac{\alpha t}{1+t^2}e^{-\alpha ^2t^2}\mathrm{d}t}\leqslant \int_0^{+\infty}{\frac{1}{\sqrt{2e}}}\cdot \frac{1}{1+t^2}\mathrm{d}t=\frac{\pi}{2\sqrt{2e}}.
\end{align*}
根据\hyperref[theorem:含参量反常积分一致收敛判别法]{Weierstrass判别法}可知该含参广义积分在 \(\alpha \in (0, +\infty)\) 上一致收敛.
\end{proof}

\begin{example}
\begin{enumerate}[(1)]
\item 讨论下列函数在 \((0,1)\) 上的连续性:
\begin{align*}
f(\alpha) = \int_{0}^{+\infty} \frac{e^{-t}}{|\sin t|^{\alpha}} \mathrm{d}t;
\end{align*}
\item 讨论下列函数在 \((0,1)\) 上的连续性:
\begin{align*}
g(\alpha) = \int_{0}^{1} \frac{f(t)}{\sqrt{|t-\alpha|}} \mathrm{d}t, \quad \text{其中 } f(t) \text{ 是 } [0,1] \text{ 上的有界可积函数}.
\end{align*}
\end{enumerate}
\end{example}
\begin{solution}
\begin{enumerate}[(1)]
\item 任取闭区间 \([\alpha_0, \beta_0] \subset (0,1)\)。对任意 \(\alpha \in [\alpha_0, \beta_0]\)，将含参量积分 \(f(\alpha)\) 按周期拆分
\begin{align}
f(\alpha) = \int_{0}^{+\infty} \frac{e^{-t}}{|\sin t|^{\alpha}} \mathrm{d}t = \sum_{n=0}^{\infty} \int_{n\pi}^{(n+1)\pi} \frac{e^{-t}}{|\sin t|^{\alpha}} \mathrm{d}t \xlongequal{u = t - n\pi} \sum_{n=0}^{\infty} e^{-n\pi} \int_{0}^{\pi} \frac{e^{-u}}{|\sin u|^{\alpha}} \mathrm{d}u. \label{eq:series_decomp_123}
\end{align}
令 \(F(\alpha) = \int_{0}^{\pi} \frac{e^{-u}}{|\sin u|^{\alpha}} \mathrm{d}u\)，则 \(f(\alpha) = \sum_{n=0}^{\infty} e^{-n\pi} F(\alpha)\)。下面先证明 \(F(\alpha)\) 在 \([\alpha_0, \beta_0]\) 上连续。
对于 \(F(\alpha)\)，我们有
\begin{align}\label{eq::jiqjw202u80u289}
F\left( \alpha \right) =\int_0^{\pi}{\frac{e^{-t}}{|\sin t|^{\alpha}}\mathrm{d}t}\leqslant \int_0^{\pi}{\frac{e^{-u}}{|\sin u|^{\beta _0}}\mathrm{d}u}\triangleq I_0.
\end{align}
下面证明控制函数 \(G(t) = \frac{e^{-t}}{|\sin t|^{\beta_0}}\) 在 \([0,+\infty)\) 上可积。将积分区间按周期拆分得
\begin{align*}
\int_0^{+\infty}{\frac{e^{-t}}{|\sin t|^{\beta _0}}\mathrm{d}t=\sum_{n=0}^{\infty}{\int_{n\pi}^{(n+1)\pi}{\frac{e^{-t}}{|\sin t|^{\beta _0}}\mathrm{d}t}}}\xlongequal{u=t-n\pi}\sum_{n=0}^{\infty}{\int_0^{\pi}{\frac{e^{-(u+n\pi )}}{|\sin u|^{\beta _0}}\mathrm{d}u}}=\sum_{n=0}^{\infty}{e^{-n\pi}\int_0^{\pi}{\frac{e^{-u}}{|\sin u|^{\beta _0}}\mathrm{d}u.}}
\end{align*}
现在证明\(I_0 = \int_{0}^{\pi} \frac{e^{-u}}{|\sin u|^{\beta_0}} \mathrm{d}u\)收敛.
由于 \(\beta_0 < 1\)，
取 \(\delta \in (0, \frac{\pi}{2})\)，当 \(u \in (0, \delta)\) 时，有 \(\sin u > \frac{1}{2} u\) 成立.从而
\begin{align*}
\int_{0}^{\delta} \frac{e^{-u}}{|\sin u|^{\beta_0}} \mathrm{d}u 
\leqslant \int_{0}^{\delta} \frac{1}{(\frac{1}{2} u)^{\beta_0}} \mathrm{d}u = 2^{\beta_0} \int_{0}^{\delta} u^{-\beta_0} \mathrm{d}u < +\infty.
\end{align*}
同理有
\begin{align*}
\int_{\pi -\delta}^{\pi}{\frac{e^{-u}}{\left| \sin u \right|^{\beta _0}}\mathrm{d}u}\xlongequal{v=\pi -u}\int_0^{\delta}{\frac{e^{v-\pi}}{\left| \sin v \right|^{\beta _0}}\mathrm{d}v}\leqslant \int_0^{\delta}{\frac{1}{\left( \frac{1}{2}v \right) ^{\beta _0}}\mathrm{d}u}<+\infty .
\end{align*}
而
\begin{align*}
\int_{\delta}^{\pi -\delta}{\frac{e^{-u}}{\left| \sin u \right|^{\beta _0}}\mathrm{d}u}\leqslant \frac{\pi -2\delta}{\left| \sin \delta \right|^{\beta _0}}<+\infty .
\end{align*}
因此
\begin{align*}
I_0=\int_0^{\delta}{\frac{e^{-u}}{\left| \sin u \right|^{\beta _0}}\mathrm{d}u}+\int_{\delta}^{\pi -\delta}{\frac{e^{-u}}{\left| \sin u \right|^{\beta _0}}\mathrm{d}u}+\int_{\pi -\delta}^{\pi}{\frac{e^{-u}}{\left| \sin u \right|^{\beta _0}}\mathrm{d}u}<+\infty .
\end{align*}
再根据\eqref{eq::jiqjw202u80u289}式和\hyperref[theorem:含参量反常积分一致收敛判别法]{Weierstrass判别法}知,\(F(\alpha)\) 在 \([\alpha_0, \beta_0]\) 上一致收敛.又因为被积函数在 \((0, \pi) \times [\alpha_0, \beta_0]\) 上连续,由\hyperref[theorem:含参量反常积分的连续性]{含参量积分连续性定理}知，\(F(\alpha)\) 在 \([\alpha_0, \beta_0]\) 上连续.

由于 \(F(\alpha)\) 在紧集 \([\alpha_0, \beta_0]\) 上连续，故存在常数 \(M > 0\)，使得对所有 \(\alpha \in [\alpha_0, \beta_0]\) 均有 \(|F(\alpha)| \leqslant M\)。于是由\eqref{eq:series_decomp_123}式可得
\begin{align*}
f\left( \alpha \right) =\sum_{n=1}^{\infty}{e^{-n\pi}F\left( \alpha \right)}\leqslant M\sum_{n=1}^{\infty}{e^{-n\pi}}<+\infty .
\end{align*}
由于几何级数 \(\sum_{n=0}^{\infty} M e^{-n\pi}\) 收敛以及\hyperref[theorem:含参量反常积分一致收敛判别法]{Weierstrass判别法}可知，函数项级数 \(f(\alpha)\) 在 \([\alpha_0, \beta_0]\) 上一致收敛。
又由于 \(F(\alpha)\) 在 \([\alpha_0, \beta_0]\) 上连续，因此每一项 \(e^{-n\pi} F(\alpha)\) 都连续，因此根据\hyperref[theorem:函数项级数连续性定理]{一致收敛级数的和函数连续性定理}，其和函数 \(f(\alpha)\) 在 \([\alpha_0, \beta_0]\) 上连续.
最后,由区间 \([\alpha_0, \beta_0] \subset (0,1)\) 的任意性可知，函数 \(f(\alpha)\) 在 \((0,1)\) 上连续。

\item 设 \(|f(t)| \leqslant M\)。任取 \(\alpha, \beta \in (0,1)\)，我们有
\begin{align*}
|g(\alpha) - g(\beta)| &= \left| \int_{0}^{1} \frac{f(t)}{\sqrt{|t-\alpha|}} \mathrm{d}t - \int_{0}^{1} \frac{f(t)}{\sqrt{|t-\beta|}} \mathrm{d}t \right| \\
&\leqslant \sup_{(0,1)} |f| \cdot \int_{0}^{1} \left| \frac{1}{\sqrt{|t-\alpha|}} - \frac{1}{\sqrt{|t-\beta|}} \right| \mathrm{d}t \\
&\leqslant \sup_{(0,1)} |f| \cdot \int_{-\infty}^{+\infty} \left| \frac{1}{\sqrt{|t-\alpha|}} - \frac{1}{\sqrt{|t-\beta|}} \right| \mathrm{d}t.
\end{align*}
令 \(t = \alpha + (\alpha - \beta)s\)，则 \(\mathrm{d}t = |\alpha - \beta| \mathrm{d}s\)，且 \(t-\alpha = (\alpha - \beta)s\)，\(t-\beta = (\alpha - \beta)(s+1)\)。代入上式得
\begin{align*}
|g(\alpha) - g(\beta)| \leqslant \sup_{(0,1)} |f| \cdot |\alpha-\beta| \int_{-\infty}^{+\infty} \frac{1}{\sqrt{|\alpha-\beta|}} \left| \frac{1}{\sqrt{|s|}} - \frac{1}{\sqrt{|s+1|}} \right| \mathrm{d}s = \sqrt{|\alpha-\beta|} \cdot \sup_{(0,1)} |f| \cdot I,
\end{align*}
其中 \(I = \int_{-\infty}^{+\infty} \left| \frac{1}{\sqrt{|s|}} - \frac{1}{\sqrt{|s+1|}} \right| \mathrm{d}s\)。下面详细计算积分 \(I\)。
\begin{align*}
I &= \int_{-\infty}^{-1} \left| \frac{1}{\sqrt{-s}} - \frac{1}{\sqrt{-(s+1)}} \right| \mathrm{d}s + \int_{-1}^{0} \left| \frac{1}{\sqrt{-s}} + \frac{1}{\sqrt{s+1}} \right| \mathrm{d}s + \int_{0}^{+\infty} \left| \frac{1}{\sqrt{s}} - \frac{1}{\sqrt{s+1}} \right| \mathrm{d}s.
\end{align*}
分别计算这三个积分。对于第一部分，令 \(u = -s\)，则
\begin{align*}
\int_{-\infty}^{-1} \left( \frac{1}{\sqrt{-s}} - \frac{1}{\sqrt{-(s+1)}} \right) \mathrm{d}s &= \int_{1}^{+\infty} \left( \frac{1}{\sqrt{u}} - \frac{1}{\sqrt{u-1}} \right) \mathrm{d}u = \lim_{M \to +\infty} \left[ 2\sqrt{u} - 2\sqrt{u-1} \right]_{1}^{M} \\
&= \lim_{M \to +\infty} \frac{2}{\sqrt{M} + \sqrt{M-1}} - (2 - 0) = -2.
\end{align*}
故其绝对值为 \(2\)。第二部分：
\begin{align*}
\int_{-1}^{0} \left( \frac{1}{\sqrt{-s}} + \frac{1}{\sqrt{s+1}} \right) \mathrm{d}s &= \left[ -2\sqrt{-s} \right]_{-1}^{0} + \left[ 2\sqrt{s+1} \right]_{-1}^{0} = (0 - (-2)) + (2 - 0) = 4.
\end{align*}
第三部分：
\begin{align*}
\int_{0}^{+\infty} \left( \frac{1}{\sqrt{s}} - \frac{1}{\sqrt{s+1}} \right) \mathrm{d}s = \lim_{M \to +\infty} \left[ 2\sqrt{s} - 2\sqrt{s+1} \right]_{0}^{M} = \lim_{M \to +\infty} \frac{-2}{\sqrt{M} + \sqrt{M+1}} - (0 - 2) = 2.
\end{align*}
由此可得 \(I = 2 + 4 + 2 = 8\)。因此
\begin{align*}
|g(\alpha) - g(\beta)| \leqslant 8 \sup_{(0,1)} |f| \cdot \sqrt{|\alpha - \beta|}.
\end{align*}
当 \(\alpha \to \beta\) 时，\(\sqrt{|\alpha - \beta|} \to 0\)，故差值趋于 \(0\)，从而 \(g(\alpha)\) 在 \((0,1)\) 上连续,且一致连续。
\end{enumerate}
\end{solution}

\subsubsection{含参量反常积分的性质}

\begin{theorem}[含参反常积分的连续性]\label{theorem:含参量反常积分的连续性}
设 \( f(x,y) \) 在 \( I \times [c, +\infty) \) 上连续，若含参量反常积分
\begin{align}\label{12}
\Phi(x) = \int_{c}^{+\infty} f(x,y) \, \mathrm{d}y
\end{align}
在 \( I \) 上(内闭)一致收敛，则 \( \Phi(x) \) 在 \( I \) 上连续。
\end{theorem}
\begin{note}
这个定理也表明，在一致收敛的条件下，极限运算与无穷积分运算可以交换：
\begin{align}\label{eq::::--34802jrw4f3454t}
\lim_{x \to x_0} \int_{c}^{+\infty} f(x,y) \, \mathrm{d}y = \int_{c}^{+\infty} f(x_0,y) \, \mathrm{d}y = \int_{c}^{+\infty} \lim_{x \to x_0} f(x,y) \, \mathrm{d}y.
\end{align}
\end{note}
\begin{proof}
由\refthe{theorem:数学分析--定理19.9}，对任一递增且趋于 \( +\infty \) 的数列 \( \{A_n\} \)（\( A_1 = c \)），函数项级数
\begin{align}\label{eq::::--34802jrw4f34t}
\Phi(x) = \sum_{n=1}^{\infty} \int_{A_n}^{A_{n+1}} f(x,y) \, \mathrm{d}y = \sum_{n=1}^{\infty} u_n(x)
\end{align}
在 \( I \) 上一致收敛。又由于 \( f(x,y) \) 在 \( I \times [c, +\infty) \) 上连续，故每个 \( u_n(x) \) 都在 \( I \) 上连续。根据函数项级数的连续性定理，函数 \( \Phi(x) \) 在 \( I \) 上连续。
\end{proof}

\begin{corollary}
设 \( f(x,y) \) 在 \( I \times [c, +\infty) \) 上连续，若 \( \Phi(x) = \int_{c}^{+\infty} f(x,y) \, \mathrm{d}y \) 在 \( I \) 上内闭一致收敛，则 \( \Phi(x) \) 在 \( I \) 上连续。
\end{corollary}

\begin{theorem}[含参反常积分的可微性]\label{theorem:含参反常积分的可微性}
设 \( f(x,y) \) 与 \( f_x(x,y) \) 在区域 \( I \times [c, +\infty) \) 上连续。若 \( \Phi(x) = \int_{c}^{+\infty} f(x,y) \, \mathrm{d}y \) 在 \( I \) 上收敛，\( \int_{c}^{+\infty} f_x(x,y) \, \mathrm{d}y \) 在 \( I \) 上内闭一致收敛，则 \( \Phi(x) \) 在 \( I \) 上可微，且
\begin{align}\label{eq::::--34802jrw4f345634t}
\Phi'(x) = \int_{c}^{+\infty} f_x(x,y) \, \mathrm{d}y.
\end{align}
\end{theorem}
\begin{note}
最后结果表明在定理条件下，求导运算和无穷积分运算可以交换。
\end{note}
\begin{proof}
对任一递增且趋于 \( +\infty \) 的数列 \( \{A_n\} \)（\( A_1 = c \)），令
\[
u_n(x) = \int_{A_n}^{A_{n+1}} f(x,y) \, \mathrm{d}y.
\]
由\refthe{theorem:含参量积分的可微性2}推得
\[
u_n'(x) = \int_{A_n}^{A_{n+1}} f_x(x,y) \, \mathrm{d}y.
\]
由 \( \int_{c}^{+\infty} f_x(x,y) \, \mathrm{d}y \) 在 \( I \) 上一致收敛及\refthe{theorem:数学分析--定理19.9}，可得函数项级数
\[
\sum_{n=1}^{\infty} u_n'(x) = \sum_{n=1}^{\infty} \int_{A_n}^{A_{n+1}} f_x(x,y) \, \mathrm{d}y
\]
在 \( I \) 上一致收敛，因此根据函数项级数的逐项求导定理，即得
\[
\Phi'(x) = \sum_{n=1}^{\infty} u_n'(x) = \sum_{n=1}^{\infty} \int_{A_n}^{A_{n+1}} f_x(x,y) \, \mathrm{d}y = \int_{c}^{+\infty} f_x(x,y) \, \mathrm{d}y,
\]
或写作
\[
\frac{\mathrm{d}}{\mathrm{d}x} \int_{c}^{+\infty} f(x,y) \, \mathrm{d}y = \int_{c}^{+\infty} \frac{\partial}{\partial x} f(x,y) \, \mathrm{d}y.
\]
\end{proof}

\begin{theorem}[可积性]\label{theorem:含参量反常积分可积性定理}
设 \( f(x,y) \) 在 \( [a,b] \times [c, +\infty) \) 上连续，若 \( \Phi(x) = \int_{c}^{+\infty} f(x,y) \, \mathrm{d}y \) 在 \( [a,b] \) 上内闭一致收敛，则 \( \Phi(x) \) 在 \( [a,b] \) 上可积，且
\begin{align}\label{eq::::--wefw234fsd3480t34t32jrwdr24t948543tdhe634}
\int_{a}^{b} \mathrm{d}x \int_{c}^{+\infty} f(x,y) \, \mathrm{d}y = \int_{c}^{+\infty} \mathrm{d}y \int_{a}^{b} f(x,y) \, \mathrm{d}x.
\end{align}
\end{theorem}
\begin{proof}
由\hyperref[theorem:含参量反常积分的连续性]{含参量反常积分的连续性}知道 \( \Phi(x) \) 在 \( [a,b] \) 上连续，从而 \( \Phi(x) \) 在 \( [a,b] \) 上可积。

又由\hyperref[theorem:含参量反常积分的连续性]{含参量反常积分的连续性}的证明中可以看到，函数项级数 \eqref{eq::::--34802jrw4f34t} 在 \( [a,b] \) 上一致收敛，且各项 \( u_n(x) \) 在 \( [a,b] \) 上连续，因此根据函数项级数逐项求积定理，有
\begin{align}\label{eq::::--wefw234fsd3480t34t32jrwdr24t94854}
\int_{a}^{b} \Phi(x) \, \mathrm{d}x = \sum_{n=1}^{\infty} \int_{a}^{b} u_n(x) \, \mathrm{d}x = \sum_{n=1}^{\infty} \int_{a}^{b} \mathrm{d}x \int_{A_n}^{A_{n+1}} f(x,y) \, \mathrm{d}y = \sum_{n=1}^{\infty} \int_{A_n}^{A_{n+1}} \mathrm{d}y \int_{a}^{b} f(x,y) \, \mathrm{d}x,
\end{align}
这里最后一步是根据\refthe{theorem:含参量正常积分的可积性定理}关于积分顺序的可交换性。\eqref{eq::::--wefw234fsd3480t34t32jrwdr24t94854} 式又可写作
\[
\int_{a}^{b} \Phi(x) \, \mathrm{d}x = \int_{c}^{+\infty} \mathrm{d}y \int_{a}^{b} f(x,y) \, \mathrm{d}x.
\]
这就是\eqref{eq::::--wefw234fsd3480t34t32jrwdr24t948543tdhe634}式。
\end{proof}

\begin{theorem}\label{theorem:含参积分可积性定理}
设 \( f(x,y) \) 在 \( [a, +\infty) \times [c, +\infty) \) 上连续。若
\begin{enumerate}[(i)]
\item \( \int_{a}^{+\infty} f(x,y) \, \mathrm{d}x \) 关于 \( y \) 在 \( [c, +\infty) \) 上内闭一致收敛，\( \int_{c}^{+\infty} f(x,y) \, \mathrm{d}y \) 关于 \( x \) 在 \( [a, +\infty) \) 上内闭一致收敛。

\item 积分
\begin{align}\label{eq::::--wefw234fsd3480t34t32jrwdr24t9485344}
\int_{a}^{+\infty} \mathrm{d}x \int_{c}^{+\infty} |f(x,y)| \, \mathrm{d}y \text{ 与 } \int_{c}^{+\infty} \mathrm{d}y \int_{a}^{+\infty} |f(x,y)| \, \mathrm{d}x
\end{align}
中有一个收敛。
\end{enumerate}
则
\begin{align}\label{eq::::--wefw234fsd3480t34t32jrwdr24t948534t34}
\int_{a}^{+\infty} \mathrm{d}x \int_{c}^{+\infty} f(x,y) \, \mathrm{d}y = \int_{c}^{+\infty} \mathrm{d}y \int_{a}^{+\infty} f(x,y) \, \mathrm{d}x.
\end{align}
\end{theorem}
\begin{proof}
不妨设 \eqref{eq::::--wefw234fsd3480t34t32jrwdr24t9485344} 式中第一个积分收敛，由此推得
\[
\int_{a}^{+\infty} \mathrm{d}x \int_{c}^{+\infty} f(x,y) \, \mathrm{d}y
\]
也收敛。当 \( d > c \) 时，
\begin{align*}
J_d &= \left| \int_{c}^{d} \mathrm{d}y \int_{a}^{+\infty} f(x,y) \, \mathrm{d}x - \int_{a}^{+\infty} \mathrm{d}x \int_{c}^{+\infty} f(x,y) \, \mathrm{d}y \right| \\
&= \left| \int_{c}^{d} \mathrm{d}y \int_{a}^{+\infty} f(x,y) \, \mathrm{d}x - \int_{a}^{+\infty} \mathrm{d}x \int_{c}^{d} f(x,y) \, \mathrm{d}y - \int_{a}^{+\infty} \mathrm{d}x \int_{d}^{+\infty} f(x,y) \, \mathrm{d}y \right|.
\end{align*}
根据条件（i）及\refthe{theorem:含参量反常积分可积性定理}，可推得
\begin{align}\label{20}
J_d = \left| \int_{a}^{+\infty} \mathrm{d}x \int_{d}^{+\infty} f(x,y) \, \mathrm{d}y \right| \leqslant \left| \int_{a}^{A} \mathrm{d}x \int_{d}^{+\infty} f(x,y) \, \mathrm{d}y \right| + \int_{A}^{+\infty} \mathrm{d}x \int_{d}^{+\infty} |f(x,y)| \, \mathrm{d}y.
\end{align}
由条件（ii），对于任给的 \( \varepsilon > 0 \)，有 \( G > a \)，使当 \( A > G \) 时，有
\[
\int_{A}^{+\infty} \mathrm{d}x \int_{d}^{+\infty} |f(x,y)| \, \mathrm{d}y < \frac{\varepsilon}{2}.
\]
选定 \( A \) 后，由 \( \int_{c}^{+\infty} f(x,y) \, \mathrm{d}y \) 的内闭一致收敛性，存在 \( M > c \)，使得当 \( d > M \) 时有
\[
\left| \int_{d}^{+\infty} f(x,y) \, \mathrm{d}y \right| < \frac{\varepsilon}{2(A - a)}, \quad \forall x \in [a,A].
\]
把这两个结果应用到 \eqref{20} 式，得到
\[
J_d < \frac{\varepsilon}{2} + \frac{\varepsilon}{2} = \varepsilon,
\]
即 \( \lim_{d \to +\infty} J_d = 0 \)，这就证明了 \eqref{eq::::--wefw234fsd3480t34t32jrwdr24t948534t34} 式。
\end{proof}

\begin{example}
记
\begin{align*}
F(y) = \int_{0}^{+\infty} \frac{\sin(xy)}{x(x^2 + 1)} \mathrm{d}x, \quad y \in \mathbb{R}.
\end{align*}
\begin{enumerate}[(1)]
\item 证明：奇函数 $F(y)$ 在 $(0,+\infty)$ 内二次可微。
\item 求出 $F(y)$ 的解析表达式。
\end{enumerate}
\end{example}
\begin{solution}
\begin{enumerate}[(1)]
\item 注意到对$\forall y\in\mathbb{R}$，有
\begin{align}
|F(y)|\leqslant \int_{0}^{+\infty} \left| \frac{\sin(xy)}{x(x^2+1)} \right| \mathrm{d}x
\leqslant \int_{0}^{+\infty} \frac{|xy|}{x(x^2+1)} \mathrm{d}x
= |y| \int_{0}^{+\infty} \frac{1}{x^2+1} \mathrm{d}x
= \frac{\pi}{2} |y| ,\label{eq::FHJIOWJFIOJJWIO32483UUUII}
\end{align}
故$F(y)$在$\mathbb{R}$上绝对收敛，进而点态收敛。又因为对$\forall y\in(0,+\infty)$，有
\begin{align*}
\int_{0}^{+\infty} \left| \frac{\partial}{\partial y} \frac{\sin(xy)}{x(x^2+1)} \right| \mathrm{d}x
= \int_{0}^{+\infty} \left| \frac{\cos(xy)}{x^2+1} \right| \mathrm{d}x
\leqslant \int_{0}^{+\infty} \frac{1}{x^2+1} \mathrm{d}x
= \frac{\pi}{2} < +\infty,
\end{align*}
所以$\int_{0}^{+\infty} \frac{\partial}{\partial y} \frac{\sin(xy)}{x(x^2+1)} \mathrm{d}x$在$(0,+\infty)$上一致收敛。故由\hyperref[theorem:含参反常积分的可微性]{含参反常积分的可微性}知$F(y)$在$(0,+\infty)$上可微，且
\begin{align}
F'(y) = \int_{0}^{+\infty} \frac{\partial}{\partial y} \frac{\sin(xy)}{x(x^2+1)} \mathrm{d}x
= \int_{0}^{+\infty} \frac{\cos(xy)}{x^2+1} \mathrm{d}x, \quad \forall y\in(0,+\infty).
\label{eq:3122187748IOUFI23--23aaaaa-1}
\end{align}
注意到
\begin{align}
\int_{0}^{+\infty} \frac{\partial}{\partial y} \frac{\cos(xy)}{x^2+1} \mathrm{d}x
= -\int_{0}^{+\infty} \frac{x\sin(xy)}{x^2+1} \mathrm{d}x
= -\int_{0}^{1} \frac{x\sin(xy)}{x^2+1} \mathrm{d}x
- \int_{1}^{+\infty} \frac{x\sin(xy)}{x^2+1} \mathrm{d}x, \quad \forall y\in(0,+\infty).
\label{eq:3122187748IOUFI23--23aaaaa-2}
\end{align}
对$\forall y\in(0,+\infty)$，有
\begin{align*}
\left| \int_{0}^{1} \frac{x\sin(xy)}{x^2+1} \mathrm{d}x \right|
\leqslant \int_{0}^{1} \left| \frac{x\sin(xy)}{x^2+1} \right| \mathrm{d}x
\leqslant \int_{0}^{1} \frac{x}{x^2+1} \mathrm{d}x
= \frac{\ln 2}{2} < +\infty.
\end{align*}
故$\int_{0}^{1} \frac{x\sin(xy)}{x^2+1} \mathrm{d}x$在$(0,+\infty)$上一致收敛。任取$[a,b] \subseteq (0,+\infty)$，则对$\forall y\in[a,b], A>0$，有
\begin{align*}
\int_{1}^{A} \sin(xy) \mathrm{d}x
= \frac{\cos y - \cos(Ay)}{y}
\leqslant \frac{2}{y}
\leqslant \frac{2}{a}.
\end{align*}
又因为$\frac{x}{x^2+1}$在$(1,+\infty)$上单调递减且$\lim\limits_{x\rightarrow+\infty} \frac{x}{x^2+1} = 0$，所以由Dirichlet判别法知$\int_{1}^{+\infty} \frac{x\sin(xy)}{x^2+1} \mathrm{d}x$在$[a,b]$上一致收敛，进而$\int_{1}^{+\infty} \frac{x\sin(xy)}{x^2+1} \mathrm{d}x$在$(0,+\infty)$上内闭一致收敛。再由\eqref{eq:3122187748IOUFI23--23aaaaa-2}式知$\int_{0}^{+\infty} \frac{\partial}{\partial y} \frac{\cos(xy)}{x^2+1} \mathrm{d}x$在$(0,+\infty)$上内闭一致收敛。从而由\hyperref[theorem:含参反常积分的可微性]{含参反常积分的可微性}知$F'(y)$在$(0,+\infty)$上可微，即$F(y)$在$(0,+\infty)$内二次可微,且
\begin{align}\label{eq::::HIUUUUUUU342424523}
F'' \left( y \right) =\int_0^{+\infty}{\frac{\partial}{\partial y}\frac{\cos \left( xy \right)}{x^2+1}\mathrm{d}x}=-\int_0^{+\infty}{\frac{x\sin \left( xy \right)}{x^2+1}\mathrm{d}x},\quad \forall y\in \left( 0,+\infty \right) .
\end{align}

\item 由\eqref{eq:3122187748IOUFI23--23aaaaa-1}
\eqref{eq::::HIUUUUUUU342424523}式知$F\in D^2(0,+\infty)$，且
\begin{align*}
F'(y)=\int_0^{+\infty}\frac{\cos(xy)}{x^2+1}\mathrm{d}x,\quad F''(y)=-\int_0^{+\infty}\frac{x\sin(xy)}{x^2+1}\mathrm{d}x,\quad \forall y\in(0,+\infty).
\end{align*}
于是对$\forall y\in(0,+\infty)$，结合\hyperref[proposition:重要定积分结论-2]{重要定积分结论\ref{proposition:重要定积分结论-2}}可得
\begin{align}
F''(y)=\int_0^{+\infty}\frac{\sin(xy)}{x(x^2+1)}\mathrm{d}x-\int_0^{+\infty}\frac{\sin(xy)}{x}\mathrm{d}x\xlongequal{t=xy}F(y)-\int_0^{+\infty}\frac{\sin t}{t}\mathrm{d}t=F(y)-\frac{\pi}{2}.\label{eq_31_0}
\end{align}
由上述微分方程解得对$\forall y\in(0,+\infty)$，有
\begin{align}
F(y)=C_1e^y+C_2e^{-y}+\frac{\pi}{2},\text{其中}C_1,C_2\text{为常数}.\label{eq_31_3}
\end{align}
由\eqref{eq::FHJIOWJFIOJJWIO32483UUUII}式知$F(y)$在$\mathbb{R}$上Lipschitz连续，进而$F\in C(\mathbb{R})$。又由题设可知$F(0)=0$，故
\begin{align}
0=F(0)=\lim\limits_{y\to0^+}F(y)=\lim\limits_{y\to0^+}\left(C_1e^y+C_2e^{-y}+\frac{\pi}{2}\right)\Longrightarrow C_1+C_2+\frac{\pi}{2}=0.\label{eq_31_4}
\end{align}
注意到对$\forall \varepsilon>0,M>0,y\in(0,+\infty)$，有
\begin{align*}
\int_0^M\frac{x\sin(xy)}{x^2+1}\mathrm{d}x&=\int_0^\varepsilon\frac{x\sin(xy)}{x^2+1}\mathrm{d}x+\int_\varepsilon^M\frac{x\sin(xy)}{x^2+1}\mathrm{d}x\\
&\leqslant\int_0^\varepsilon\frac{x}{x^2+1}\mathrm{d}x+\int_\varepsilon^M\sin(xy)\mathrm{d}x\\
&\leqslant\varepsilon+\frac{\cos(\varepsilon y)-\cos(My)}{y}\leqslant\varepsilon+\frac{2}{y}.
\end{align*}
令$M,y\to+\infty$，再令$\varepsilon\to0^+$得
\begin{align*}
\lim\limits_{y\to+\infty}\int_0^{+\infty}\frac{x\sin(xy)}{x^2+1}\mathrm{d}x=0.
\end{align*}
从而结合\eqref{eq_31_0}式和\hyperref[proposition:重要定积分结论-2]{重要定积分结论\ref{proposition:重要定积分结论-2}}知
\begin{align*}
\lim\limits_{y\to+\infty}F(y)=\lim\limits_{y\to+\infty}\int_0^{+\infty}\frac{\sin(xy)}{x(x^2+1)}\mathrm{d}x=\int_0^{+\infty}\frac{\sin t}{t}\mathrm{d}t-\lim\limits_{y\to+\infty}\int_0^{+\infty}\frac{x\sin(xy)}{x^2+1}\mathrm{d}x=\frac{\pi}{2}.
\end{align*}
再利用\eqref{eq_31_3}式可得
\begin{align*}
\lim\limits_{y\to+\infty}F(y)=\lim\limits_{y\to+\infty}\left(C_1e^y+C_2e^{-y}+\frac{\pi}{2}\right)=\lim\limits_{y\to+\infty}\left(C_1e^y+\frac{\pi}{2}\right)=\frac{\pi}{2}\Longrightarrow C_1=0.
\end{align*}
再结合\eqref{eq_31_4}式可得$C_2=-\frac{\pi}{2}$。故再由\eqref{eq_31_3}式知
\begin{align*}
F(y)=-\frac{\pi}{2}e^{-y}+\frac{\pi}{2}=\frac{\pi}{2}(1-e^{-y}),\quad \forall y\in(0,+\infty).
\end{align*}
又由题设可知$F(y)$是$\mathbb{R}$上的奇函数，因此
\begin{align*}
F(y)=
\begin{cases}
\frac{\pi}{2}(1-e^{-y}),&y\in(0,+\infty),\\
0,&y=0,\\
-\frac{\pi}{2}(1-e^y),&y\in(-\infty,0).
\end{cases}
\end{align*}
此即
\begin{align*}
F(y)=\frac{\pi}{2}\sgn(y)\left(1-e^{-|y|}\right),\quad \forall y\in\mathbb{R}.
\end{align*}
\end{enumerate}
\end{solution}






\subsubsection{含参量瑕积分}

\begin{definition}
设 \( f(x,y) \) 在区域 \( R = [a,b] \times [c,d) \) 上有定义。若对 \( x \) 的某些值，\( y = d \) 为函数 \( f(x,y) \) 的瑕点，则称
\begin{align}\label{eq::::-23535232-wefw234fsd3480t34t32jrwdr24t}
\int_{c}^{d} f(x,y) \mathrm{d}y
\end{align}
为\textbf{含参量 \( \boldsymbol{x} \) 的无界函数反常积分}，或简称为含参量反常积分。若对每一个 \( x \in [a,b] \)，积分 \eqref{eq::::-23535232-wefw234fsd3480t34t32jrwdr24t} 都收敛，则其积分值是 \( x \) 在 \( [a,b] \) 上取值的函数。
\end{definition}

\begin{definition}
对任给正数 \( \varepsilon \)，总存在某正数 \( \delta < d - c \)，使得当 \( 0 < \eta < \delta \) 时，对一切 \( x \in [a,b] \)，都有
\begin{align*}
\left| \int_{d - \eta}^{d} f(x,y) \mathrm{d}y \right| < \varepsilon,
\end{align*}
则称含参量反常积分 \eqref{eq::::-23535232-wefw234fsd3480t34t32jrwdr24t} 在 \( [a,b] \) 上一致收敛。
\end{definition}















\end{document}